\section{Bursting against budding}
\label{dist:sec:budding}

This chapter opened by objecting to the constant drip. It can now say
precisely what that objection amounts to, and at the level where it was made:
one cell.

The comparator is the budding cell of the constant-rate models. In its
simplest stochastic form a budding infected cell experiences two
\emph{independent} event types --- release of one infectious particle at rate
$p$, and death of the cell at rate $d_{\Icell}$ --- and the two continue until
the cell dies. Its survival and mean release are
\begin{equation}
I_{\rm bud}(t)=\ee^{-d_{\Icell}t},\qquad
V_{\rm bud}(t)=\frac{p}{d_{\Icell}}\lrb{1-\ee^{-d_{\Icell}t}},
\label{dist:eq:bud}
\end{equation}
the same one-cell limit the Basic Model of Viral Replication reaches at the
population level. The release count of a budding cell is geometric on
$\{0,1,2,\ldots\}$:
\begin{equation}
\pr{\Kb_{\rm bud}=k}=\frac{d_{\Icell}}{d_{\Icell}+p}
\lrb{\frac{p}{d_{\Icell}+p}}^k,\qquad k=0,1,2,\ldots,
\label{dist:eq:budgeom}
\end{equation}
since each successive event is a release with probability
$p/(d_{\Icell}+p)$ and the terminating death with probability
$d_{\Icell}/(d_{\Icell}+p)$.

The structural contrast with bursting is exact: in budding, release and death
are independent events, and a cell may release many particles and then die
with nothing left to release; in bursting, release and death are the
\emph{same} event, and whatever has accumulated is released at once. That
difference is why a bursting cell's release law is a statement about its
accumulated load, and why \cref{dist:sec:burst-size} has no budding analogue:
a cell that releases continuously holds nothing to be quasi-stationary about.
\Cref{dist:tab:budburst} sets the two descriptions side by side for the
$\mu=0$ bursting case.

\begin{table}[tbp]
\centering
\renewcommand{\arraystretch}{1.35}
\small
\begin{tabular}{@{}lll@{}}
\toprule
 & budding & bursting ($\mu=0$) \\
\midrule
$\ddt{I}$ & $-d_{\Icell} I$ & $-(\lambda+\delta)I+\lambda I^2$ \\
$I(t)$ & $\ee^{-d_{\Icell}t}$ & $\dfrac{\lambda+\delta}{\lambda+\delta\,\ee^{(\lambda+\delta)t}}$ \\
mean cell lifetime & $1/d_{\Icell}$ & $\dfrac{1}{\lambda}\log\lrb{1+\dfrac{\lambda}{\delta}}$ \\
$V(t)$ & $\dfrac{p}{d_{\Icell}}(1-I)$ & $\dfrac{\lambda}{\delta}+1-\dfrac{2\lambda+\delta}{\delta}I+\dfrac{\lambda}{\delta}I^2$ \\
release law & geometric on $\{0,1,\ldots\}$, \cref{dist:eq:budgeom} & geometric on $\{1,2,\ldots\}$, $\pr{\Kb=k}=\frac{\delta}{\lambda}a^{-k}$ \\
\bottomrule
\end{tabular}
\caption{Single-cell budding and bursting compared, $\mu=0$ for bursting.
Budding parameters: release rate $p$ and infected-cell death rate
$d_{\Icell}$. Bursting parameters: division rate $\lambda$ and rupture rate
$\delta$ per intracellular particle; $a=(\lambda+\delta)/\lambda$ here.}
\label{dist:tab:budburst}
\end{table}

Two rows of the table carry the comparison. The mean cell lifetime of a
bursting cell grows only logarithmically with $\lambda/\delta$ while its mean
release grows linearly; budding offers no such coupling, its lifetime and its
release being controlled by separate parameters. And the release laws, both
geometric, differ in their support: budding can release nothing at all from a
cell that died early, while bursting never releases nothing from a cell that
ruptured. A trade-off along an axis of this kind is a familiar theme in the
evolution of virulence \cite{gilchrist2006evolution}.

Matching the means does not repair the difference. \Cref{dist:fig:budding}
sets the two release laws against each other with the budding mean tuned to
the BDC conditional mean $V_\infty/(1-b)=17.23$: the budding law still holds
probability $0.055$ of releasing nothing. Matching one moment matches neither
the support nor the tails. That is the concrete form of the objection this
chapter opened with: a fitted constant-rate model can reproduce the average
yield of a lytic cell and still be wrong about the quantity that decides
whether a small inoculum establishes at all.

The comparison is deliberately confined to one cell. What happens when these
descriptions are embedded in a population, where the timing of release governs
establishment and invasion speed, belongs to \ChPop; the pathogenesis-level
comparison belongs to
\fwd{path}{\ChPath}. All three
chapters return to this contrast, and answer different questions with it.

\begin{figure}[tbp]
\centering
\includegraphics[width=0.92\textwidth]{fig15_burst_vs_budding/N4a_3_burst_vs_budding_pmf}
\caption{Release-count laws for the BDC process at
$(\lambda,\mu,\delta)=(1,0.2,0.05)$ and a budding comparator whose mean
$p/d_{\Icell}$ is matched to the BDC conditional mean
$V_\infty/(1-b)=17.23$. Panel~(a): the BDC law conditioned on bursting is
geometric on $\{1,2,\ldots\}$, whereas the budding law is geometric on
$\{0,1,2,\ldots\}$ and retains probability $0.055$ of zero release.
Panel~(b): survival functions $\pr{\Kb\ge k}$ and
$\pr{\Kb_{\rm bud}\ge k}$. Matching the means matches neither the support
nor the tails.}
\label{dist:fig:budding}
\end{figure}
